% Options for packages loaded elsewhere
\PassOptionsToPackage{unicode}{hyperref}
\PassOptionsToPackage{hyphens}{url}
\PassOptionsToPackage{dvipsnames,svgnames,x11names}{xcolor}
%
\documentclass[
]{scrartcl}

\usepackage{amsmath,amssymb}
\usepackage{iftex}
\ifPDFTeX
  \usepackage[T1]{fontenc}
  \usepackage[utf8]{inputenc}
  \usepackage{textcomp} % provide euro and other symbols
\else % if luatex or xetex
  \usepackage{unicode-math}
  \defaultfontfeatures{Scale=MatchLowercase}
  \defaultfontfeatures[\rmfamily]{Ligatures=TeX,Scale=1}
\fi
\usepackage{lmodern}
\ifPDFTeX\else  
    % xetex/luatex font selection
\fi
% Use upquote if available, for straight quotes in verbatim environments
\IfFileExists{upquote.sty}{\usepackage{upquote}}{}
\IfFileExists{microtype.sty}{% use microtype if available
  \usepackage[]{microtype}
  \UseMicrotypeSet[protrusion]{basicmath} % disable protrusion for tt fonts
}{}
\makeatletter
\@ifundefined{KOMAClassName}{% if non-KOMA class
  \IfFileExists{parskip.sty}{%
    \usepackage{parskip}
  }{% else
    \setlength{\parindent}{0pt}
    \setlength{\parskip}{6pt plus 2pt minus 1pt}}
}{% if KOMA class
  \KOMAoptions{parskip=half}}
\makeatother
\usepackage{xcolor}
\setlength{\emergencystretch}{3em} % prevent overfull lines
\setcounter{secnumdepth}{5}
% Make \paragraph and \subparagraph free-standing
\ifx\paragraph\undefined\else
  \let\oldparagraph\paragraph
  \renewcommand{\paragraph}[1]{\oldparagraph{#1}\mbox{}}
\fi
\ifx\subparagraph\undefined\else
  \let\oldsubparagraph\subparagraph
  \renewcommand{\subparagraph}[1]{\oldsubparagraph{#1}\mbox{}}
\fi

% Lower the position of page numbers
\usepackage{fancyhdr}
\fancyhf{}
\fancyfoot[C]{\thepage}
\setlength{\footskip}{100pt} % Increase this value to lower the page numbers


\usepackage{color}
\usepackage{fancyvrb}
\newcommand{\VerbBar}{|}
\newcommand{\VERB}{\Verb[commandchars=\\\{\}]}
\DefineVerbatimEnvironment{Highlighting}{Verbatim}{commandchars=\\\{\}}
% Add ',fontsize=\small' for more characters per line
\usepackage{framed}
\definecolor{shadecolor}{RGB}{241,243,245}
\newenvironment{Shaded}{\begin{snugshade}}{\end{snugshade}}
\newcommand{\AlertTok}[1]{\textcolor[rgb]{0.68,0.00,0.00}{#1}}
\newcommand{\AnnotationTok}[1]{\textcolor[rgb]{0.37,0.37,0.37}{#1}}
\newcommand{\AttributeTok}[1]{\textcolor[rgb]{0.40,0.45,0.13}{#1}}
\newcommand{\BaseNTok}[1]{\textcolor[rgb]{0.68,0.00,0.00}{#1}}
\newcommand{\BuiltInTok}[1]{\textcolor[rgb]{0.00,0.23,0.31}{#1}}
\newcommand{\CharTok}[1]{\textcolor[rgb]{0.13,0.47,0.30}{#1}}
\newcommand{\CommentTok}[1]{\textcolor[rgb]{0.37,0.37,0.37}{#1}}
\newcommand{\CommentVarTok}[1]{\textcolor[rgb]{0.37,0.37,0.37}{\textit{#1}}}
\newcommand{\ConstantTok}[1]{\textcolor[rgb]{0.56,0.35,0.01}{#1}}
\newcommand{\ControlFlowTok}[1]{\textcolor[rgb]{0.00,0.23,0.31}{\textbf{#1}}}
\newcommand{\DataTypeTok}[1]{\textcolor[rgb]{0.68,0.00,0.00}{#1}}
\newcommand{\DecValTok}[1]{\textcolor[rgb]{0.68,0.00,0.00}{#1}}
\newcommand{\DocumentationTok}[1]{\textcolor[rgb]{0.37,0.37,0.37}{\textit{#1}}}
\newcommand{\ErrorTok}[1]{\textcolor[rgb]{0.68,0.00,0.00}{#1}}
\newcommand{\ExtensionTok}[1]{\textcolor[rgb]{0.00,0.23,0.31}{#1}}
\newcommand{\FloatTok}[1]{\textcolor[rgb]{0.68,0.00,0.00}{#1}}
\newcommand{\FunctionTok}[1]{\textcolor[rgb]{0.28,0.35,0.67}{#1}}
\newcommand{\ImportTok}[1]{\textcolor[rgb]{0.00,0.46,0.62}{#1}}
\newcommand{\InformationTok}[1]{\textcolor[rgb]{0.37,0.37,0.37}{#1}}
\newcommand{\KeywordTok}[1]{\textcolor[rgb]{0.00,0.23,0.31}{\textbf{#1}}}
\newcommand{\NormalTok}[1]{\textcolor[rgb]{0.00,0.23,0.31}{#1}}
\newcommand{\OperatorTok}[1]{\textcolor[rgb]{0.37,0.37,0.37}{#1}}
\newcommand{\OtherTok}[1]{\textcolor[rgb]{0.00,0.23,0.31}{#1}}
\newcommand{\PreprocessorTok}[1]{\textcolor[rgb]{0.68,0.00,0.00}{#1}}
\newcommand{\RegionMarkerTok}[1]{\textcolor[rgb]{0.00,0.23,0.31}{#1}}
\newcommand{\SpecialCharTok}[1]{\textcolor[rgb]{0.37,0.37,0.37}{#1}}
\newcommand{\SpecialStringTok}[1]{\textcolor[rgb]{0.13,0.47,0.30}{#1}}
\newcommand{\StringTok}[1]{\textcolor[rgb]{0.13,0.47,0.30}{#1}}
\newcommand{\VariableTok}[1]{\textcolor[rgb]{0.07,0.07,0.07}{#1}}
\newcommand{\VerbatimStringTok}[1]{\textcolor[rgb]{0.13,0.47,0.30}{#1}}
\newcommand{\WarningTok}[1]{\textcolor[rgb]{0.37,0.37,0.37}{\textit{#1}}}

\providecommand{\tightlist}{%
  \setlength{\itemsep}{0pt}\setlength{\parskip}{0pt}}\usepackage{longtable,booktabs,array}
\usepackage{calc} % for calculating minipage widths
% Correct order of tables after \paragraph or \subparagraph
\usepackage{etoolbox}
\makeatletter
\patchcmd\longtable{\par}{\if@noskipsec\mbox{}\fi\par}{}{}
\makeatother
% Allow footnotes in longtable head/foot
\IfFileExists{footnotehyper.sty}{\usepackage{footnotehyper}}{\usepackage{footnote}}
\makesavenoteenv{longtable}
\usepackage{graphicx}
\makeatletter
\newsavebox\pandoc@box
\newcommand*\pandocbounded[1]{% scales image to fit in text height/width
  \sbox\pandoc@box{#1}%
  \Gscale@div\@tempa{\textheight}{\dimexpr\ht\pandoc@box+\dp\pandoc@box\relax}%
  \Gscale@div\@tempb{\linewidth}{\wd\pandoc@box}%
  \ifdim\@tempb\p@<\@tempa\p@\let\@tempa\@tempb\fi% select the smaller of both
  \ifdim\@tempa\p@<\p@\scalebox{\@tempa}{\usebox\pandoc@box}%
  \else\usebox{\pandoc@box}%
  \fi%
}
% Set default figure placement to htbp
\def\fps@figure{htbp}
\makeatother

% Change font size
\usepackage[font={footnotesize, sf}, labelfont=bf, margin=0.05\linewidth]{caption}
\renewcommand{\familydefault}{\sfdefault}
% Bold the 'Figure #' in the caption and separate it from the title/caption
% with a period
% Captions will be left justified
\usepackage[
	aboveskip=1pt,
	labelfont=bf,
	labelsep=period,
	justification=raggedright,
	singlelinecheck=off
]{caption}
\renewcommand{\figurename}{Fig}

% Add numbered lines
\usepackage{lineno}
% \linenumbers

% Package to include multiple title pages
% This will allow me to add a tile to the main text and to the SI
\usepackage{titling}

% Define command to begin the supplementary section
\newcommand{\beginsupplement}{
\setcounter{page}{1} % Restart page counter
\setcounter{section}{0} % Restart section counter
% \renewcommand{\thesection}{\Alph{section}}%
\renewcommand{\thesection}{}
\renewcommand{\thesubsection}{\Alph{subsection}}
\setcounter{table}{0} % Restart table counter
\renewcommand{\thetable}{\Alph{table}}
\setcounter{figure}{0} % Restart figure counter
\renewcommand{\thefigure}{\Alph{figure}}%
\setcounter{equation}{0} % Restart equation counter
\renewcommand{\theequation}{S\arabic{equation}}%
}

% Add author affiliations
\usepackage{authblk}

% Allow to define personalized colors
\usepackage{xcolor}
\makeatletter
\@ifpackageloaded{caption}{}{\usepackage{caption}}
\AtBeginDocument{%
\ifdefined\contentsname
  \renewcommand*\contentsname{Table of contents}
\else
  \newcommand\contentsname{Table of contents}
\fi
\ifdefined\listfigurename
  \renewcommand*\listfigurename{List of Figures}
\else
  \newcommand\listfigurename{List of Figures}
\fi
\ifdefined\listtablename
  \renewcommand*\listtablename{List of Tables}
\else
  \newcommand\listtablename{List of Tables}
\fi
\ifdefined\figurename
  \renewcommand*\figurename{Figure}
\else
  \newcommand\figurename{Figure}
\fi
\ifdefined\tablename
  \renewcommand*\tablename{Table}
\else
  \newcommand\tablename{Table}
\fi
}
\@ifpackageloaded{float}{}{\usepackage{float}}
\floatstyle{ruled}
\@ifundefined{c@chapter}{\newfloat{codelisting}{h}{lop}}{\newfloat{codelisting}{h}{lop}[chapter]}
\floatname{codelisting}{Listing}
\newcommand*\listoflistings{\listof{codelisting}{List of Listings}}
\makeatother
\makeatletter
\makeatother
\makeatletter
\@ifpackageloaded{caption}{}{\usepackage{caption}}
\@ifpackageloaded{subcaption}{}{\usepackage{subcaption}}
\makeatother
\ifLuaTeX
  \usepackage{selnolig}  % disable illegal ligatures
\fi

%%%%%%%%%%%%%%%%%%%%%%%%%%%%%%%%%%%%%%%%%%%%%%%%%%%%%%%%%%%%%%%%%%%%%%%%%%%%%%%%
\usepackage[
    style=vancouver,
    sorting=none,				% Do not sort bibliography
	url=false, 					% Do not show url in reference
	doi=false, 					% Do not show doi in reference
	isbn=false, 				% Do not show isbn link in reference
	eprint=false, 			    % Do not show eprint link in reference
	maxbibnames=9, 			    % Include up to 9 names in citation
	firstinits=true,
]{biblatex}
%%%%%%%%%%%%%%%%%%%%%%%%%%%%%%%%%%%%%%%%%%%%%%%%%%%%%%%%%%%%%%%%%%%%%%%%%%%%%%%%
\addbibresource{../references.bib}
\IfFileExists{bookmark.sty}{\usepackage{bookmark}}{\usepackage{hyperref}}
\IfFileExists{xurl.sty}{\usepackage{xurl}}{} % add URL line breaks if available
\urlstyle{same} % disable monospaced font for URLs
\hypersetup{
  pdftitle={Evolutionary Dynamics Simulation},
  pdfauthor={Manuel Razo-Mejia},
  pdfkeywords={evolutionary dynamics, fitness
landscapes, genotype-phenotype density, Metropolis-Hastings evolutionary
algorithm},
  colorlinks=true,
  linkcolor={blue},
  filecolor={Maroon},
  citecolor={Blue},
  urlcolor={Blue},
  pdfcreator={LaTeX via pandoc}}

%%%%%%%%%%%%%%%%%%%%%%%%%%%%%%%%%%%%%%%%%%%%%%%%%%%%%%%%%%%%%%%%%%%%%%%%%%%%%%%%
% Add title
\title{Evolutionary Dynamics Simulation}
%%%%%%%%%%%%%%%%%%%%%%%%%%%%%%%%%%%%%%%%%%%%%%%%%%%%%%%%%%%%%%%%%%%%%%%%%%%%%%%%

%%%%%%%%%%%%%%%%%%%%%%%%%%%%%%%%%%%%%%%%%%%%%%%%%%%%%%%%%%%%%%%%%%%%%%%%%%%%%%%%
% Add list of authors

% Loop through authors
\author[1]{Manuel Razo-Mejia}

% Loop through affiliations
\affil[1]{Department of Biology, Stanford University}

% Add corresponding author
%%%%%%%%%%%%%%%%%%%%%%%%%%%%%%%%%%%%%%%%%%%%%%%%%%%%%%%%%%%%%%%%%%%%%%%%%%%%%%%%

%%%%%%%%%%%%%%%%%%%%%%%%%%%%%%%%%%%%%%%%%%%%%%%%%%%%%%%%%%%%%%%%%%%%%%%%%%%%%%%%
% Set affiliations in small font
\renewcommand\Affilfont{\itshape\small}
%%%%%%%%%%%%%%%%%%%%%%%%%%%%%%%%%%%%%%%%%%%%%%%%%%%%%%%%%%%%%%%%%%%%%%%%%%%%%%%%

%%%%%%%%%%%%%%%%%%%%%%%%%%%%%%%%%%%%%%%%%%%%%%%%%%%%%%%%%%%%%%%%%%%%%%%%%%%%%%%%
% Remove date
\date{}
%%%%%%%%%%%%%%%%%%%%%%%%%%%%%%%%%%%%%%%%%%%%%%%%%%%%%%%%%%%%%%%%%%%%%%%%%%%%%%%%

% Add package for line numbering
\usepackage{lineno}

\begin{document}
\linenumbers
\maketitle
(c) This work is licensed under a
\href{https://creativecommons.org/licenses/by/4.0/}{Creative Commons
Attribution License CC-BY 4.0}. All code contained herein is licensed
under an \href{https://opensource.org/licenses/MIT}{MIT license}.

\section{Evolutionary Dynamics
Simulation}\label{evolutionary-dynamics-simulation}

In this notebook, we simulate the dynamics of populations evolving in a
2-dimensional phenotypic space under various fitness landscapes and a
fixed mutational landscape. We'll use the Metropolis-Hastings algorithm
to model evolution as a biased random walk influenced by both fitness
and mutational constraints.

\subsection{Setup Environment}\label{setup-environment}

First, let's import the necessary packages for our simulation.

\begin{Shaded}
\begin{Highlighting}[]
\CommentTok{\# Import custom project package}
\ImportTok{import} \BuiltInTok{Antibiotic}
\ImportTok{import} \BuiltInTok{Antibiotic.mh}\NormalTok{ as mh }\CommentTok{\# Metropolis{-}Hastings dynamics module}

\CommentTok{\# Import packages for storing results}
\ImportTok{import} \BuiltInTok{DimensionalData}\NormalTok{ as DD}

\CommentTok{\# Import JLD2 for saving results}
\ImportTok{import} \BuiltInTok{JLD2}

\CommentTok{\# Import IterTools for iterating over Cartesian products}
\ImportTok{import} \BuiltInTok{IterTools}

\CommentTok{\# Import basic math libraries}
\ImportTok{import} \BuiltInTok{StatsBase}
\ImportTok{import} \BuiltInTok{LinearAlgebra}
\ImportTok{import} \BuiltInTok{Random}
\ImportTok{import} \BuiltInTok{Distributions}
\ImportTok{import} \BuiltInTok{Distances}

\CommentTok{\# Packages for visualization}
\ImportTok{using} \BuiltInTok{CairoMakie}
\ImportTok{import} \BuiltInTok{ColorSchemes}

\CommentTok{\# Activate Makie backend}
\NormalTok{CairoMakie.}\FunctionTok{activate!}\NormalTok{()}

\CommentTok{\# Set custom plotting style}
\NormalTok{Antibiotic.viz.}\FunctionTok{theme\_makie!}\NormalTok{()}

\CommentTok{\# Set random seed for reproducibility}
\BuiltInTok{Random}\NormalTok{.}\FunctionTok{seed!}\NormalTok{(}\FloatTok{42}\NormalTok{)}
\end{Highlighting}
\end{Shaded}

\subsection{Simulation Parameters}\label{simulation-parameters}

Let's define the parameters that govern our evolutionary simulation. The
things we need to define are:

\begin{itemize}
\tightlist
\item
  The dimensionality of the phenotype space \texttt{n\_dim}
\item
  The number of phenotypic initial conditions, i.e., number of lineages
  (positions in phenotype space) \texttt{n\_sim}
\item
  The number of replicates (evolving strains per lineage)
  \texttt{n\_rep}
\item
  The inverse temperature (controls selection strength) \texttt{β}
\item
  The mutation step size \texttt{µ}
\item
  The number of evolution steps \texttt{n\_steps}
\item
  The number of fitness landscapes \texttt{n\_fit\_lans}
\item
  The range of peak means \texttt{peak\_mean\_min} and
  \texttt{peak\_mean\_max}
\item
  The range of fitness amplitudes \texttt{fit\_amp\_min} and
  \texttt{fit\_amp\_max}
\item
  The range of fitness covariances \texttt{fit\_cov\_min} and
  \texttt{fit\_cov\_max}
\item
  The number of fitness peaks \texttt{n\_fit\_peaks\_min} and
  \texttt{n\_fit\_peaks\_max}
\end{itemize}

\begin{Shaded}
\begin{Highlighting}[]
\CommentTok{\# Phenotype space dimensionality}
\NormalTok{n\_dim }\OperatorTok{=} \FloatTok{2}
\CommentTok{\# Number of initial conditions (positions in phenotype space)}
\NormalTok{n\_sim }\OperatorTok{=} \FloatTok{10}
\CommentTok{\# Number of replicates (evolving strains per initial condition)}
\NormalTok{n\_rep }\OperatorTok{=} \FloatTok{2}
\CommentTok{\# Inverse temperature (controls selection strength)}
\NormalTok{β }\OperatorTok{=} \FloatTok{10.0}
\CommentTok{\# Mutation step size}
\NormalTok{µ }\OperatorTok{=} \FloatTok{0.1}
\CommentTok{\# Number of evolution steps}
\NormalTok{n\_steps }\OperatorTok{=} \FloatTok{300}

\CommentTok{\# Define number of fitness landscapes}
\NormalTok{n\_fit\_lans }\OperatorTok{=} \FloatTok{50}

\CommentTok{\# Define range of peak means}
\NormalTok{peak\_mean\_min }\OperatorTok{=} \OperatorTok{{-}}\FloatTok{4.0}
\NormalTok{peak\_mean\_max }\OperatorTok{=} \FloatTok{4.0}

\CommentTok{\# Define range of fitness amplitudes}
\NormalTok{fit\_amp\_min }\OperatorTok{=} \FloatTok{1.0}
\NormalTok{fit\_amp\_max }\OperatorTok{=} \FloatTok{5.0}

\CommentTok{\# Define covariance range}
\NormalTok{fit\_cov\_min }\OperatorTok{=} \FloatTok{3.0}
\NormalTok{fit\_cov\_max }\OperatorTok{=} \FloatTok{10.0}

\CommentTok{\# Define possible number of fitness peaks}
\NormalTok{n\_fit\_peaks\_min }\OperatorTok{=} \FloatTok{1}
\NormalTok{n\_fit\_peaks\_max }\OperatorTok{=} \FloatTok{4}
\end{Highlighting}
\end{Shaded}

Having defined these parameters, we must define the two functions that
govern the dynamics of the system:

\begin{itemize}
\tightlist
\item
  The genotype-phenotype density function \(GP(\underline{p})\)
\item
  The fitness functions \(F_E(\underline{p})\) for each environment
  \(E\).
\end{itemize}

Let's start by defining the genotype-phenotype density function.

\subsection{Genotype-Phenotype Density
Function}\label{genotype-phenotype-density-function}

The genotype-phenotype density function \(GP(\underline{p})\) represents
constraints on phenotypic evolution due to the underlying
genotype-phenotype mapping. It defines which regions of phenotypic space
are more or less accessible to mutations.

In our simulation, we model the genotype-phenotype density function as a
collection of inverted Gaussian peaks. These peaks define areas of
phenotypic space with lower densities of possible mutations. In other
words, we define four ``wells'' in which the genotype-phenotype density
drops, making it more unlikely for mutations to take a population
towards these regions.

For illustrative purposes, as done in the main text, we will define the
genotype-phenotype density function as a collection of four Gaussian
peaks, which we will place at the corners of a square in phenotype
space.

\begin{Shaded}
\begin{Highlighting}[]
\CommentTok{\# GP density peak amplitude}
\NormalTok{gp\_evo\_amplitude }\OperatorTok{=} \FloatTok{1.0}

\CommentTok{\# GP density peak means {-} defines 4 peaks in a square arrangement}
\NormalTok{gp\_means }\OperatorTok{=}\NormalTok{ [}
\NormalTok{    [}\OperatorTok{{-}}\FloatTok{1.5}\NormalTok{, }\OperatorTok{{-}}\FloatTok{1.5}\NormalTok{],}
\NormalTok{    [}\FloatTok{1.5}\NormalTok{, }\OperatorTok{{-}}\FloatTok{1.5}\NormalTok{],}
\NormalTok{    [}\FloatTok{1.5}\NormalTok{, }\FloatTok{1.5}\NormalTok{],}
\NormalTok{    [}\OperatorTok{{-}}\FloatTok{1.5}\NormalTok{, }\FloatTok{1.5}\NormalTok{],}
\NormalTok{]}

\CommentTok{\# GP density peak covariance (controls spread of each peak)}
\NormalTok{gp\_evo\_covariance }\OperatorTok{=} \FloatTok{0.45}

\CommentTok{\# Create GP density peaks object using the \textasciigrave{}GaussianPeaks\textasciigrave{} constructor}
\NormalTok{gp\_evo\_peaks }\OperatorTok{=}\NormalTok{ mh.}\FunctionTok{GaussianPeaks}\NormalTok{(}
\NormalTok{    gp\_evo\_amplitude,}
\NormalTok{    gp\_means,}
\NormalTok{    gp\_evo\_covariance}
\NormalTok{)}

\CommentTok{\# Define grid on which to evaluate GP density landscape}
\NormalTok{gp\_evo\_grid }\OperatorTok{=} \FunctionTok{range}\NormalTok{(peak\_mean\_min, peak\_mean\_max, length}\OperatorTok{=}\FloatTok{100}\NormalTok{)}

\CommentTok{\# Evaluate GP density landscape on grid}
\NormalTok{gp\_evo\_grid\_points }\OperatorTok{=}\NormalTok{ mh.}\FunctionTok{genetic\_density}\NormalTok{(}
    \FunctionTok{tuple}\NormalTok{(}\FunctionTok{repeat}\NormalTok{([gp\_evo\_grid], n\_dim)}\OperatorTok{...}\NormalTok{),}
\NormalTok{    gp\_evo\_peaks}
\NormalTok{)}

\CommentTok{\# Define grid of possible initial conditions}
\NormalTok{init\_grid }\OperatorTok{=}\NormalTok{ [[x}\OperatorTok{...}\NormalTok{] for x }\KeywordTok{in}\NormalTok{ IterTools.}\FunctionTok{product}\NormalTok{(}\FunctionTok{fill}\NormalTok{(gp\_evo\_grid, }\FloatTok{2}\NormalTok{)}\OperatorTok{...}\NormalTok{)]}
\end{Highlighting}
\end{Shaded}

Let's visualize this mutational landscape:

\begin{Shaded}
\begin{Highlighting}[]
\CommentTok{\# Define ranges of phenotypes to evaluate}
\NormalTok{x }\OperatorTok{=} \FunctionTok{range}\NormalTok{(}\OperatorTok{{-}}\FloatTok{4}\NormalTok{, }\FloatTok{4}\NormalTok{, length}\OperatorTok{=}\FloatTok{100}\NormalTok{)}
\NormalTok{y }\OperatorTok{=} \FunctionTok{range}\NormalTok{(}\OperatorTok{{-}}\FloatTok{4}\NormalTok{, }\FloatTok{4}\NormalTok{, length}\OperatorTok{=}\FloatTok{100}\NormalTok{)}

\CommentTok{\# Initialize figure}
\NormalTok{fig }\OperatorTok{=} \FunctionTok{Figure}\NormalTok{(size}\OperatorTok{=}\NormalTok{(}\FloatTok{350}\NormalTok{, }\FloatTok{300}\NormalTok{))}

\CommentTok{\# Add axis}
\NormalTok{ax }\OperatorTok{=} \FunctionTok{Axis}\NormalTok{(}
\NormalTok{    fig[}\FloatTok{1}\NormalTok{, }\FloatTok{1}\NormalTok{],}
\NormalTok{    xlabel}\OperatorTok{=}\StringTok{"phenotype 1"}\NormalTok{,}
\NormalTok{    ylabel}\OperatorTok{=}\StringTok{"phenotype 2"}\NormalTok{,}
\NormalTok{    aspect}\OperatorTok{=}\FunctionTok{AxisAspect}\NormalTok{(}\FloatTok{1}\NormalTok{),}
\NormalTok{    title}\OperatorTok{=}\StringTok{"Genotype{-}to{-}Phenotype}\SpecialCharTok{\textbackslash{}n}\StringTok{Density"}\NormalTok{,}
\NormalTok{    titlesize}\OperatorTok{=}\FloatTok{14}\NormalTok{,}
\NormalTok{    yticklabelsvisible}\OperatorTok{=}\ConstantTok{false}\NormalTok{,}
\NormalTok{    xticklabelsvisible}\OperatorTok{=}\ConstantTok{false}\NormalTok{,}
\NormalTok{)}

\CommentTok{\# Evaluate mutational landscape}
\NormalTok{M }\OperatorTok{=}\NormalTok{ mh.}\FunctionTok{genetic\_density}\NormalTok{(x, y, gp\_evo\_peaks)}

\CommentTok{\# Plot mutational landscape}
\NormalTok{hm\_gp }\OperatorTok{=} \FunctionTok{heatmap!}\NormalTok{(ax, x, y, M, colormap}\OperatorTok{=}\FunctionTok{Reverse}\NormalTok{(ColorSchemes.Purples\_9))}

\CommentTok{\# Add contour plot}
\FunctionTok{contour!}\NormalTok{(ax, x, y, M, color}\OperatorTok{=:}\NormalTok{white)}

\CommentTok{\# Add colorbar for mutational landscape}
\FunctionTok{Colorbar}\NormalTok{(}
\NormalTok{    fig[}\FloatTok{1}\NormalTok{, }\FloatTok{2}\NormalTok{],}
\NormalTok{    hm\_gp,}
\NormalTok{    label}\OperatorTok{=}\StringTok{"genotype{-}phenotype}\SpecialCharTok{\textbackslash{}n}\StringTok{density"}\NormalTok{,}
\NormalTok{    labelsize}\OperatorTok{=}\FloatTok{13}\NormalTok{,}
\NormalTok{    height}\OperatorTok{=}\FunctionTok{Relative}\NormalTok{(}\FloatTok{1}\NormalTok{),}
\NormalTok{    ticksvisible}\OperatorTok{=}\ConstantTok{false}\NormalTok{,}
\NormalTok{    ticklabelsvisible}\OperatorTok{=}\ConstantTok{false}
\NormalTok{)}

\NormalTok{fig}
\end{Highlighting}
\end{Shaded}

\includegraphics[width=7.29167in,height=6.25in]{evolutionary_dynamics_files/figure-pdf/cell-5-output-1.png}

The mutational landscape shows four peaks (in darker purple)
representing regions of phenotypic space that have lower accessibility
through mutations. The lighter regions between these peaks represent
``valleys'' which are much easier for populations to traverse.

Let's continue by defining the fitness landscapes that we'll use to
simulate evolutionary dynamics.

\subsection{Fixed Evolution Condition}\label{fixed-evolution-condition}

We first define a reference fitness landscape (evolution condition) that
we'll use as a baseline for comparison. This is the simplest possible
fitness landscape with a single peak in the center of the phenotype
space.

\begin{Shaded}
\begin{Highlighting}[]
\CommentTok{\# Evolution condition amplitude}
\NormalTok{fit\_evo\_amplitude }\OperatorTok{=} \FloatTok{5.0}
\CommentTok{\# Evolution condition mean}
\NormalTok{fit\_evo\_mean }\OperatorTok{=}\NormalTok{ [}\FloatTok{0.0}\NormalTok{, }\FloatTok{0.0}\NormalTok{]}
\CommentTok{\# Evolution condition covariance}
\NormalTok{fit\_evo\_covariance }\OperatorTok{=} \FloatTok{3.0}
\CommentTok{\# Create peak}
\NormalTok{fit\_evo\_peak }\OperatorTok{=}\NormalTok{ mh.}\FunctionTok{GaussianPeak}\NormalTok{(}
\NormalTok{    fit\_evo\_amplitude,}
\NormalTok{    fit\_evo\_mean,}
\NormalTok{    fit\_evo\_covariance}
\NormalTok{)}
\end{Highlighting}
\end{Shaded}

Let's visualize this reference fitness landscape:

\begin{Shaded}
\begin{Highlighting}[]
\CommentTok{\# Initialize figure}
\NormalTok{fig }\OperatorTok{=} \FunctionTok{Figure}\NormalTok{(size}\OperatorTok{=}\NormalTok{(}\FloatTok{350}\NormalTok{, }\FloatTok{300}\NormalTok{))}

\CommentTok{\# Add axis}
\NormalTok{ax }\OperatorTok{=} \FunctionTok{Axis}\NormalTok{(}
\NormalTok{    fig[}\FloatTok{1}\NormalTok{, }\FloatTok{1}\NormalTok{],}
\NormalTok{    xlabel}\OperatorTok{=}\StringTok{"phenotype 1"}\NormalTok{,}
\NormalTok{    ylabel}\OperatorTok{=}\StringTok{"phenotype 2"}\NormalTok{,}
\NormalTok{    aspect}\OperatorTok{=}\FunctionTok{AxisAspect}\NormalTok{(}\FloatTok{1}\NormalTok{),}
\NormalTok{    title}\OperatorTok{=}\StringTok{"reference fitness landscape"}\NormalTok{,}
\NormalTok{    titlesize}\OperatorTok{=}\FloatTok{14}\NormalTok{,}
\NormalTok{    yticklabelsvisible}\OperatorTok{=}\ConstantTok{false}\NormalTok{,}
\NormalTok{    xticklabelsvisible}\OperatorTok{=}\ConstantTok{false}\NormalTok{,}
\NormalTok{)}

\CommentTok{\# Evaluate fitness landscape}
\NormalTok{F }\OperatorTok{=}\NormalTok{ mh.}\FunctionTok{fitness}\NormalTok{(x, y, fit\_evo\_peak)}

\CommentTok{\# Plot fitness landscape}
\NormalTok{hm\_fit }\OperatorTok{=} \FunctionTok{heatmap!}\NormalTok{(ax, x, y, F, colormap}\OperatorTok{=:}\NormalTok{algae)}

\CommentTok{\# Add contour plot}
\FunctionTok{contour!}\NormalTok{(ax, x, y, F, color}\OperatorTok{=:}\NormalTok{white)}

\CommentTok{\# Add colorbar for fitness landscape}
\FunctionTok{Colorbar}\NormalTok{(}
\NormalTok{    fig[}\FloatTok{1}\NormalTok{, }\FloatTok{2}\NormalTok{],}
\NormalTok{    hm\_fit,}
\NormalTok{    label}\OperatorTok{=}\StringTok{"fitness"}\NormalTok{,}
\NormalTok{    labelsize}\OperatorTok{=}\FloatTok{13}\NormalTok{,}
\NormalTok{    height}\OperatorTok{=}\FunctionTok{Relative}\NormalTok{(}\FloatTok{1}\NormalTok{),}
\NormalTok{    ticksvisible}\OperatorTok{=}\ConstantTok{false}\NormalTok{,}
\NormalTok{    ticklabelsvisible}\OperatorTok{=}\ConstantTok{false}
\NormalTok{)}

\NormalTok{fig}
\end{Highlighting}
\end{Shaded}

\includegraphics[width=7.29167in,height=6.25in]{evolutionary_dynamics_files/figure-pdf/cell-7-output-1.png}

This fitness landscape shows a single peak centered at
\texttt{{[}0,\ 0{]}}. The color represents fitness, with darker regions
having higher fitness. The contour lines show equal fitness values.

\subsection{Alternative Fitness
Landscapes}\label{alternative-fitness-landscapes}

To explore how evolution proceeds in different environments, we generate
a variety of fitness landscapes with randomly selected numbers,
positions, and shapes of fitness peaks. The sampled numbers of fitness
peaks, positions, and amplitudes will be drawn from uniform
distributions within the ranges defined above.

\begin{Shaded}
\begin{Highlighting}[]
\CommentTok{\# Reset random seed for reproducibility}
\BuiltInTok{Random}\NormalTok{.}\FunctionTok{seed!}\NormalTok{(}\FloatTok{42}\NormalTok{)}

\CommentTok{\# Initialize array to hold fitness landscapes}
\NormalTok{fit\_lans }\OperatorTok{=} \FunctionTok{Array}\DataTypeTok{\{mh.AbstractPeak\}}\NormalTok{(}\ConstantTok{undef}\NormalTok{, n\_fit\_lans)}

\CommentTok{\# Add fixed evolution condition to fitness landscapes}
\NormalTok{fit\_lans[}\FloatTok{1}\NormalTok{] }\OperatorTok{=}\NormalTok{ fit\_evo\_peak}

\CommentTok{\# Loop over number of fitness landscapes}
\ControlFlowTok{for}\NormalTok{ i }\KeywordTok{in} \FloatTok{2}\OperatorTok{:}\NormalTok{n\_fit\_lans}
    \CommentTok{\# Sample number of fitness peaks}
\NormalTok{    n\_fit\_peaks }\OperatorTok{=} \FunctionTok{rand}\NormalTok{(n\_fit\_peaks\_min}\OperatorTok{:}\NormalTok{n\_fit\_peaks\_max)}

    \CommentTok{\# Sample fitness means as 2D vectors from uniform distribution}
\NormalTok{    fit\_means }\OperatorTok{=}\NormalTok{ [}
        \FunctionTok{rand}\NormalTok{(Distributions.}\FunctionTok{Uniform}\NormalTok{(peak\_mean\_min, peak\_mean\_max), }\FloatTok{2}\NormalTok{)}
\NormalTok{        for \_ }\KeywordTok{in} \FloatTok{1}\OperatorTok{:}\NormalTok{n\_fit\_peaks}
\NormalTok{    ]}

    \CommentTok{\# Sample fitness amplitudes from uniform distribution}
\NormalTok{    fit\_amplitudes }\OperatorTok{=} \FunctionTok{rand}\NormalTok{(}
\NormalTok{        Distributions.}\FunctionTok{Uniform}\NormalTok{(fit\_amp\_min, fit\_amp\_max), n\_fit\_peaks}
\NormalTok{    )}

    \CommentTok{\# Sample fitness covariances from uniform distribution}
\NormalTok{    fit\_covariances }\OperatorTok{=} \FunctionTok{rand}\NormalTok{(}
\NormalTok{        Distributions.}\FunctionTok{Uniform}\NormalTok{(fit\_cov\_min, fit\_cov\_max), n\_fit\_peaks}
\NormalTok{    )}

    \CommentTok{\# Check dimensionality}
    \ControlFlowTok{if}\NormalTok{ n\_fit\_peaks }\OperatorTok{==} \FloatTok{1}
        \CommentTok{\# Create fitness peaks}
\NormalTok{        fit\_lans[i] }\OperatorTok{=}\NormalTok{ mh.}\FunctionTok{GaussianPeak}\NormalTok{(}
            \FunctionTok{first}\NormalTok{(fit\_amplitudes), }\FunctionTok{first}\NormalTok{(fit\_means), }\FunctionTok{first}\NormalTok{(fit\_covariances)}
\NormalTok{        )}
    \ControlFlowTok{else}
        \CommentTok{\# Create fitness peaks}
\NormalTok{        fit\_lans[i] }\OperatorTok{=}\NormalTok{ mh.}\FunctionTok{GaussianPeaks}\NormalTok{(}
\NormalTok{            fit\_amplitudes, fit\_means, fit\_covariances}
\NormalTok{        )}
    \ControlFlowTok{end}
\ControlFlowTok{end}
\end{Highlighting}
\end{Shaded}

Let's visualize a sample of these fitness landscapes:

\begin{Shaded}
\begin{Highlighting}[]
\BuiltInTok{Random}\NormalTok{.}\FunctionTok{seed!}\NormalTok{(}\FloatTok{42}\NormalTok{)}

\CommentTok{\# Define number of rows and columns}
\NormalTok{n\_rows }\OperatorTok{=} \FloatTok{4}
\NormalTok{n\_cols }\OperatorTok{=} \FloatTok{4}

\CommentTok{\# Define number of landscapes to display}
\NormalTok{n\_sample }\OperatorTok{=}\NormalTok{ n\_rows }\OperatorTok{*}\NormalTok{ n\_cols}

\CommentTok{\# Initialize figure}
\NormalTok{fig }\OperatorTok{=} \FunctionTok{Figure}\NormalTok{(size}\OperatorTok{=}\NormalTok{(}\FloatTok{600}\NormalTok{, }\FloatTok{600}\NormalTok{))}

\CommentTok{\# Add global grid layout}
\NormalTok{gl }\OperatorTok{=}\NormalTok{ fig[}\FloatTok{1}\NormalTok{, }\FloatTok{1}\NormalTok{] }\OperatorTok{=} \FunctionTok{GridLayout}\NormalTok{()}

\CommentTok{\# Loop over a sample of fitness landscapes}
\ControlFlowTok{for}\NormalTok{ (i, idx) }\KeywordTok{in} \FunctionTok{enumerate}\NormalTok{(}\FloatTok{1}\OperatorTok{:}\NormalTok{n\_sample)}
    \CommentTok{\# Define row and column indices}
\NormalTok{    row }\OperatorTok{=}\NormalTok{ (i }\OperatorTok{{-}} \FloatTok{1}\NormalTok{) }\OperatorTok{÷}\NormalTok{ n\_cols }\OperatorTok{+} \FloatTok{1}
\NormalTok{    col }\OperatorTok{=}\NormalTok{ (i }\OperatorTok{{-}} \FloatTok{1}\NormalTok{) }\OperatorTok{\%}\NormalTok{ n\_cols }\OperatorTok{+} \FloatTok{1}

    \CommentTok{\# Create subplot}
\NormalTok{    ax }\OperatorTok{=} \FunctionTok{Axis}\NormalTok{(}
\NormalTok{        gl[row, col],}
\NormalTok{        aspect}\OperatorTok{=}\FunctionTok{AxisAspect}\NormalTok{(}\FloatTok{1}\NormalTok{),}
\NormalTok{        title}\OperatorTok{=}\StringTok{"landscape }\SpecialCharTok{$}\NormalTok{(idx)}\StringTok{"}\NormalTok{,}
\NormalTok{        titlesize}\OperatorTok{=}\FloatTok{12}\NormalTok{,}
\NormalTok{        yticklabelsvisible}\OperatorTok{=}\ConstantTok{false}\NormalTok{,}
\NormalTok{        xticklabelsvisible}\OperatorTok{=}\ConstantTok{false}\NormalTok{,}
\NormalTok{    )}

    \CommentTok{\# Evaluate fitness landscape}
\NormalTok{    F }\OperatorTok{=}\NormalTok{ mh.}\FunctionTok{fitness}\NormalTok{(x, y, fit\_lans[idx])}

    \CommentTok{\# Plot fitness landscape}
\NormalTok{    hm\_fit }\OperatorTok{=} \FunctionTok{heatmap!}\NormalTok{(ax, x, y, F, colormap}\OperatorTok{=:}\NormalTok{algae)}

    \CommentTok{\# Add contour plot}
    \FunctionTok{contour!}\NormalTok{(ax, x, y, F, color}\OperatorTok{=:}\NormalTok{white)}
\ControlFlowTok{end}

\CommentTok{\# Set global axis labels}
\FunctionTok{Label}\NormalTok{(}
\NormalTok{    fig[}\OperatorTok{:}\NormalTok{, }\KeywordTok{end}\NormalTok{, }\FunctionTok{Bottom}\NormalTok{()],}
    \StringTok{"phenotype 1"}\NormalTok{,}
\NormalTok{    fontsize}\OperatorTok{=}\FloatTok{14}\NormalTok{,}
\NormalTok{    padding}\OperatorTok{=}\NormalTok{(}\FloatTok{0}\NormalTok{, }\FloatTok{0}\NormalTok{, }\FloatTok{0}\NormalTok{, }\FloatTok{10}\NormalTok{),}
\NormalTok{)}

\FunctionTok{Label}\NormalTok{(}
\NormalTok{    fig[}\FloatTok{1}\NormalTok{, }\OperatorTok{:}\NormalTok{, }\FunctionTok{Left}\NormalTok{()],}
    \StringTok{"phenotype 2"}\NormalTok{,}
\NormalTok{    fontsize}\OperatorTok{=}\FloatTok{14}\NormalTok{,}
\NormalTok{    rotation}\OperatorTok{=}\ConstantTok{π} \OperatorTok{/} \FloatTok{2}\NormalTok{,}
\NormalTok{)}

\NormalTok{fig}
\end{Highlighting}
\end{Shaded}

\includegraphics[width=12.5in,height=12.5in]{evolutionary_dynamics_files/figure-pdf/cell-9-output-1.png}

These fitness landscapes represent different environments in which our
populations will evolve. Each landscape has a unique configuration of
fitness peaks, with varying heights, positions, and shapes.

We are now ready to simulate the evolutionary dynamics of our
populations.

\subsection{Simulating Evolution with
Metropolis-Hastings}\label{simulating-evolution-with-metropolis-hastings}

Now we'll simulate the evolutionary dynamics using the
Metropolis-Hastings algorithm. In this algorithm:

\begin{enumerate}
\def\labelenumi{\arabic{enumi}.}
\tightlist
\item
  We start with an initial phenotype
\item
  We propose a random mutation (a step in phenotype space)
\item
  We accept or reject the mutation based on both fitness and
  genotype-phenotype density
\item
  We repeat for many generations
\end{enumerate}

The Metropolis-Hastings acceptance probability is given by:

\begin{equation}\phantomsection\label{eq-acceptance-probability}{
P_{\text{accept}} = \min\left(
    1,
    \frac{F_E(\underline{x}')^β \cdot GP(\underline{x}')^β}
    {F_E(\underline{x})^β \cdot GP(\underline{x})^β}
\right)
}\end{equation}

where \(F_E(\underline{x})\) is the fitness in environment \(E\),
\(GP(\underline{x})\) is the genotype-phenotype density, and \(\beta\)
is the inverse temperature parameter controlling selection strength.

Our first step is to sample initial conditions for our replicates. For
this, we sample initial positions on phenotype space from a uniform
distribution, taking into account the mutational landscape (not to start
at a low mutational peak).

\begin{Shaded}
\begin{Highlighting}[]
\CommentTok{\# Reset random seed for reproducibility}
\BuiltInTok{Random}\NormalTok{.}\FunctionTok{seed!}\NormalTok{(}\FloatTok{42}\NormalTok{)}

\CommentTok{\# Sample initial positions on phenotype space from uniform distribution taking}
\CommentTok{\# into account mutational landscape (not to start at a low mutational peak)}
\NormalTok{x0 }\OperatorTok{=}\NormalTok{ StatsBase.}\FunctionTok{sample}\NormalTok{(}
    \FunctionTok{vec}\NormalTok{(init\_grid),}
\NormalTok{    StatsBase.}\FunctionTok{Weights}\NormalTok{(}\FunctionTok{vec}\NormalTok{(gp\_evo\_grid\_points)),}
\NormalTok{    n\_sim}
\NormalTok{)}

\CommentTok{\# Select initial conditions for replicates from multivariate normal distribution}
\CommentTok{\# around each initial condition}
\NormalTok{x0\_reps }\OperatorTok{=} \FunctionTok{reduce}\NormalTok{(}
\NormalTok{    (x, y) }\OperatorTok{{-}\textgreater{}} \FunctionTok{cat}\NormalTok{(x, y, dims}\OperatorTok{=}\FloatTok{3}\NormalTok{),}
\NormalTok{    [}
        \FunctionTok{rand}\NormalTok{(Distributions.}\FunctionTok{MvNormal}\NormalTok{(x0[i], }\FloatTok{0.1}\NormalTok{), n\_rep)}
\NormalTok{        for i }\KeywordTok{in} \FloatTok{1}\OperatorTok{:}\NormalTok{n\_sim}
\NormalTok{    ]}
\NormalTok{)}
\CommentTok{\# Change order of dimensions}
\NormalTok{x0\_reps }\OperatorTok{=} \FunctionTok{permutedims}\NormalTok{(x0\_reps, (}\FloatTok{1}\NormalTok{, }\FloatTok{3}\NormalTok{, }\FloatTok{2}\NormalTok{))}
\end{Highlighting}
\end{Shaded}

Next, we initialize a very convenient data structure from the
\href{https://rafaqz.github.io/DimensionalData.jl/v0.29.15/}{DimensionalData.jl}
package to hold our trajectories and fitness values for each replicate,
landscape, and evolution condition.

\begin{Shaded}
\begin{Highlighting}[]
\CommentTok{\# Define dimensions to be used with DimensionalData}
\NormalTok{phenotype }\OperatorTok{=}\NormalTok{ DD.}\FunctionTok{Dim}\DataTypeTok{\{:phenotype\}}\NormalTok{([}\OperatorTok{:}\NormalTok{x1, }\OperatorTok{:}\NormalTok{x2]) }\CommentTok{\# phenotype}
\NormalTok{fitness }\OperatorTok{=}\NormalTok{ DD.}\FunctionTok{Dim}\DataTypeTok{\{:fitness\}}\NormalTok{([}\OperatorTok{:}\NormalTok{fitness]) }\CommentTok{\# fitness}
\NormalTok{time }\OperatorTok{=}\NormalTok{ DD.}\FunctionTok{Dim}\DataTypeTok{\{:time\}}\NormalTok{(}\FloatTok{0}\OperatorTok{:}\NormalTok{n\_steps) }\CommentTok{\# time}
\NormalTok{lineage }\OperatorTok{=}\NormalTok{ DD.}\FunctionTok{Dim}\DataTypeTok{\{:lineage\}}\NormalTok{(}\FloatTok{1}\OperatorTok{:}\NormalTok{n\_sim) }\CommentTok{\# lineage}
\NormalTok{replicate }\OperatorTok{=}\NormalTok{ DD.}\FunctionTok{Dim}\DataTypeTok{\{:replicate\}}\NormalTok{(}\FloatTok{1}\OperatorTok{:}\NormalTok{n\_rep) }\CommentTok{\# replicate}
\NormalTok{landscape }\OperatorTok{=}\NormalTok{ DD.}\FunctionTok{Dim}\DataTypeTok{\{:landscape\}}\NormalTok{(}\FloatTok{1}\OperatorTok{:}\NormalTok{n\_fit\_lans) }\CommentTok{\# landscape}
\NormalTok{evo }\OperatorTok{=}\NormalTok{ DD.}\FunctionTok{Dim}\DataTypeTok{\{:evo\}}\NormalTok{(}\FloatTok{1}\OperatorTok{:}\NormalTok{n\_fit\_lans) }\CommentTok{\# evolution condition}

\CommentTok{\# Initialize DimensionalData array to hold trajectories and fitness}
\NormalTok{phenotype\_traj }\OperatorTok{=}\NormalTok{ DD.}\FunctionTok{zeros}\NormalTok{(}
    \DataTypeTok{Float32}\NormalTok{,}
\NormalTok{    phenotype,}
\NormalTok{    time,}
\NormalTok{    lineage,}
\NormalTok{    replicate,}
\NormalTok{    landscape,}
\NormalTok{    evo,}
\NormalTok{)}
\NormalTok{fitness\_traj }\OperatorTok{=}\NormalTok{ DD.}\FunctionTok{zeros}\NormalTok{(}
    \DataTypeTok{Float32}\NormalTok{,}
\NormalTok{    fitness,}
\NormalTok{    time,}
\NormalTok{    lineage,}
\NormalTok{    replicate,}
\NormalTok{    landscape,}
\NormalTok{    evo,}
\NormalTok{)}

\CommentTok{\# Stack arrays to store trajectories in phenotype and fitness dimensions}
\NormalTok{x\_traj }\OperatorTok{=}\NormalTok{ DD.}\FunctionTok{DimStack}\NormalTok{(}
\NormalTok{    (phenotype}\OperatorTok{=}\NormalTok{phenotype\_traj, fitness}\OperatorTok{=}\NormalTok{fitness\_traj),}
\NormalTok{)}

\CommentTok{\# Store initial conditions}
\NormalTok{x\_traj.phenotype[time}\OperatorTok{=}\FloatTok{1}\NormalTok{] }\OperatorTok{=} \FunctionTok{repeat}\NormalTok{(}
\NormalTok{    x0\_reps,}
\NormalTok{    outer}\OperatorTok{=}\NormalTok{(}\FloatTok{1}\NormalTok{, }\FloatTok{1}\NormalTok{, }\FloatTok{1}\NormalTok{, n\_fit\_lans, n\_fit\_lans)}
\NormalTok{)}

\CommentTok{\# Map initial phenotype to fitness}
\NormalTok{x\_traj.fitness[time}\OperatorTok{=}\FloatTok{1}\NormalTok{] }\OperatorTok{=} \FunctionTok{repeat}\NormalTok{(}
    \FunctionTok{reduce}\NormalTok{(}
\NormalTok{        (x, y) }\OperatorTok{{-}\textgreater{}} \FunctionTok{cat}\NormalTok{(x, y, dims}\OperatorTok{=}\FloatTok{3}\NormalTok{),}
\NormalTok{        mh.}\FunctionTok{fitness}\NormalTok{.(}\FunctionTok{Ref}\NormalTok{(x0\_reps), fit\_lans)}
\NormalTok{    ),}
\NormalTok{    outer}\OperatorTok{=}\NormalTok{(}\FloatTok{1}\NormalTok{, }\FloatTok{1}\NormalTok{, }\FloatTok{1}\NormalTok{, }\FloatTok{1}\NormalTok{, n\_fit\_lans)}
\NormalTok{)}
\end{Highlighting}
\end{Shaded}

Now we can loop over landscapes, lineages, and replicates to simulate
the evolutionary dynamics of our populations. The core function to run
the simulations is already implemented in the \texttt{mh} module of our
custom package.

\begin{Shaded}
\begin{Highlighting}[]
\BuiltInTok{Random}\NormalTok{.}\FunctionTok{seed!}\NormalTok{(}\FloatTok{42}\NormalTok{)}

\CommentTok{\# Loop over landscapes}
\ControlFlowTok{for}\NormalTok{ evo }\KeywordTok{in}\NormalTok{ DD.}\FunctionTok{dims}\NormalTok{(x\_traj, }\OperatorTok{:}\NormalTok{evo)}
    \CommentTok{\# Loop over lineages}
    \ControlFlowTok{for}\NormalTok{ lin }\KeywordTok{in}\NormalTok{ DD.}\FunctionTok{dims}\NormalTok{(x\_traj, }\OperatorTok{:}\NormalTok{lineage)}
        \CommentTok{\# Loop over replicates}
        \ControlFlowTok{for}\NormalTok{ rep }\KeywordTok{in}\NormalTok{ DD.}\FunctionTok{dims}\NormalTok{(x\_traj, }\OperatorTok{:}\NormalTok{replicate)}
            \CommentTok{\# Run Metropolis{-}Hastings algorithm}
\NormalTok{            trajectory }\OperatorTok{=}\NormalTok{ mh.}\FunctionTok{evo\_metropolis\_hastings}\NormalTok{(}
\NormalTok{                x\_traj.phenotype[}
\NormalTok{                    time}\OperatorTok{=}\FloatTok{1}\NormalTok{,}
\NormalTok{                    lineage}\OperatorTok{=}\NormalTok{lin,}
\NormalTok{                    replicate}\OperatorTok{=}\NormalTok{rep,}
\NormalTok{                    landscape}\OperatorTok{=}\NormalTok{evo,}
\NormalTok{                    evo}\OperatorTok{=}\NormalTok{evo,}
\NormalTok{                ].data,}
\NormalTok{                fit\_lans[evo],}
\NormalTok{                gp\_evo\_peaks,}
\NormalTok{                β,}
\NormalTok{                µ,}
\NormalTok{                n\_steps}
\NormalTok{            )}

            \CommentTok{\# Store trajectory}
\NormalTok{            x\_traj.phenotype[}
\NormalTok{                lineage}\OperatorTok{=}\NormalTok{lin,}
\NormalTok{                replicate}\OperatorTok{=}\NormalTok{rep,}
\NormalTok{                evo}\OperatorTok{=}\NormalTok{evo,}
\NormalTok{            ] }\OperatorTok{.=}\NormalTok{ trajectory}

            \CommentTok{\# Calculate and store fitness for each point in the trajectory}
\NormalTok{            x\_traj.fitness[}
\NormalTok{                lineage}\OperatorTok{=}\NormalTok{lin,}
\NormalTok{                replicate}\OperatorTok{=}\NormalTok{rep,}
\NormalTok{                evo}\OperatorTok{=}\NormalTok{evo,}
\NormalTok{            ] }\OperatorTok{=}\NormalTok{ mh.}\FunctionTok{fitness}\NormalTok{(trajectory, fit\_lans)}
        \ControlFlowTok{end}
    \ControlFlowTok{end}
\ControlFlowTok{end}
\end{Highlighting}
\end{Shaded}

\subsection{Visualize Evolutionary
Trajectories}\label{visualize-evolutionary-trajectories}

Having run the simulations, we can now visualize some of the
evolutionary trajectories to understand how populations navigate the
phenotypic space.

\begin{Shaded}
\begin{Highlighting}[]
\BuiltInTok{Random}\NormalTok{.}\FunctionTok{seed!}\NormalTok{(}\FloatTok{42}\NormalTok{)}

\CommentTok{\# Define number of rows and columns}
\NormalTok{n\_rows }\OperatorTok{=} \FloatTok{4}
\NormalTok{n\_cols }\OperatorTok{=} \FloatTok{4}

\CommentTok{\# Define number of landscapes to display}
\NormalTok{n\_sample }\OperatorTok{=}\NormalTok{ n\_rows }\OperatorTok{*}\NormalTok{ n\_cols}

\CommentTok{\# Initialize figure}
\NormalTok{fig }\OperatorTok{=} \FunctionTok{Figure}\NormalTok{(size}\OperatorTok{=}\NormalTok{(}\FloatTok{600}\NormalTok{, }\FloatTok{600}\NormalTok{))}

\CommentTok{\# Add global grid layout}
\NormalTok{gl }\OperatorTok{=}\NormalTok{ fig[}\FloatTok{1}\NormalTok{, }\FloatTok{1}\NormalTok{] }\OperatorTok{=} \FunctionTok{GridLayout}\NormalTok{()}

\CommentTok{\# Loop over a sample of fitness landscapes}
\ControlFlowTok{for}\NormalTok{ (i, idx) }\KeywordTok{in} \FunctionTok{enumerate}\NormalTok{(}\FloatTok{1}\OperatorTok{:}\NormalTok{n\_sample)}
    \CommentTok{\# Define row and column indices}
\NormalTok{    row }\OperatorTok{=}\NormalTok{ (i }\OperatorTok{{-}} \FloatTok{1}\NormalTok{) }\OperatorTok{÷}\NormalTok{ n\_cols }\OperatorTok{+} \FloatTok{1}
\NormalTok{    col }\OperatorTok{=}\NormalTok{ (i }\OperatorTok{{-}} \FloatTok{1}\NormalTok{) }\OperatorTok{\%}\NormalTok{ n\_cols }\OperatorTok{+} \FloatTok{1}

    \CommentTok{\# Create subplot}
\NormalTok{    ax }\OperatorTok{=} \FunctionTok{Axis}\NormalTok{(}
\NormalTok{        gl[row, col],}
\NormalTok{        aspect}\OperatorTok{=}\FunctionTok{AxisAspect}\NormalTok{(}\FloatTok{1}\NormalTok{),}
\NormalTok{        title}\OperatorTok{=}\StringTok{"landscape }\SpecialCharTok{$}\NormalTok{(idx)}\StringTok{"}\NormalTok{,}
\NormalTok{        titlesize}\OperatorTok{=}\FloatTok{12}\NormalTok{,}
\NormalTok{        yticklabelsvisible}\OperatorTok{=}\ConstantTok{false}\NormalTok{,}
\NormalTok{        xticklabelsvisible}\OperatorTok{=}\ConstantTok{false}\NormalTok{,}
\NormalTok{    )}

    \CommentTok{\# Evaluate fitness landscape}
\NormalTok{    F }\OperatorTok{=}\NormalTok{ mh.}\FunctionTok{fitness}\NormalTok{(x, y, fit\_lans[idx])}

    \CommentTok{\# Plot fitness landscape}
\NormalTok{    hm\_fit }\OperatorTok{=} \FunctionTok{heatmap!}\NormalTok{(ax, x, y, F, colormap}\OperatorTok{=:}\NormalTok{algae)}

    \CommentTok{\# Add contour plot}
    \FunctionTok{contour!}\NormalTok{(ax, x, y, F, color}\OperatorTok{=:}\NormalTok{black)}

    \CommentTok{\# Evaluate genetic density}
\NormalTok{    M }\OperatorTok{=}\NormalTok{ mh.}\FunctionTok{genetic\_density}\NormalTok{(x, y, gp\_evo\_peaks)}

    \CommentTok{\# Plot genetic density as contour lines}
    \FunctionTok{contour!}\NormalTok{(ax, x, y, M, color}\OperatorTok{=:}\NormalTok{white)}

    \CommentTok{\# Set limits}
    \FunctionTok{xlims!}\NormalTok{(ax, }\OperatorTok{{-}}\FloatTok{4}\NormalTok{, }\FloatTok{4}\NormalTok{)}
    \FunctionTok{ylims!}\NormalTok{(ax, }\OperatorTok{{-}}\FloatTok{4}\NormalTok{, }\FloatTok{4}\NormalTok{)}

    \CommentTok{\# Initialize counter for color index}
\NormalTok{    color\_idx }\OperatorTok{=} \FloatTok{0}

    \CommentTok{\# Plot evolutionary trajectories}
    \ControlFlowTok{for}\NormalTok{ lin }\KeywordTok{in}\NormalTok{ DD.}\FunctionTok{dims}\NormalTok{(x\_traj, }\OperatorTok{:}\NormalTok{lineage)}
        \ControlFlowTok{for}\NormalTok{ rep }\KeywordTok{in}\NormalTok{ DD.}\FunctionTok{dims}\NormalTok{(x\_traj, }\OperatorTok{:}\NormalTok{replicate)}
            \CommentTok{\# Extract x and y coordinates}
\NormalTok{            x\_data }\OperatorTok{=}\NormalTok{ x\_traj.phenotype[}
                \CommentTok{\# time=DD.At(1:n\_steps+1),}
\NormalTok{                phenotype}\OperatorTok{=}\NormalTok{DD.}\FunctionTok{At}\NormalTok{(}\OperatorTok{:}\NormalTok{x1),}
\NormalTok{                lineage}\OperatorTok{=}\NormalTok{lin,}
\NormalTok{                replicate}\OperatorTok{=}\NormalTok{rep,}
\NormalTok{                landscape}\OperatorTok{=}\NormalTok{idx,}
\NormalTok{                evo}\OperatorTok{=}\NormalTok{idx,}
\NormalTok{            ].data}
\NormalTok{            y\_data }\OperatorTok{=}\NormalTok{ x\_traj.phenotype[}
                \CommentTok{\# time=DD.At(1:n\_steps+1),}
\NormalTok{                phenotype}\OperatorTok{=}\NormalTok{DD.}\FunctionTok{At}\NormalTok{(}\OperatorTok{:}\NormalTok{x2),}
\NormalTok{                lineage}\OperatorTok{=}\NormalTok{lin,}
\NormalTok{                replicate}\OperatorTok{=}\NormalTok{rep,}
\NormalTok{                landscape}\OperatorTok{=}\NormalTok{idx,}
\NormalTok{                evo}\OperatorTok{=}\NormalTok{idx,}
\NormalTok{            ].data}

            \CommentTok{\# Plot trajectory}
\NormalTok{            color\_idx }\OperatorTok{+=} \FloatTok{1}
            \FunctionTok{lines!}\NormalTok{(}
\NormalTok{                ax,}
\NormalTok{                x\_data,}
\NormalTok{                y\_data,}
\NormalTok{                color}\OperatorTok{=}\NormalTok{ColorSchemes.glasbey\_hv\_n256[color\_idx],}
\NormalTok{                linewidth}\OperatorTok{=}\FloatTok{1}
\NormalTok{            )}

            \CommentTok{\# Plot initial condition}
            \FunctionTok{scatter!}\NormalTok{(}
\NormalTok{                ax,}
\NormalTok{                x\_data[}\FloatTok{1}\NormalTok{],}
\NormalTok{                y\_data[}\FloatTok{1}\NormalTok{],}
\NormalTok{                color}\OperatorTok{=}\NormalTok{ColorSchemes.glasbey\_hv\_n256[color\_idx],}
\NormalTok{                markersize}\OperatorTok{=}\FloatTok{12}\NormalTok{,}
\NormalTok{                marker}\OperatorTok{=:}\NormalTok{xcross,}
\NormalTok{                strokecolor}\OperatorTok{=:}\NormalTok{black,}
\NormalTok{                strokewidth}\OperatorTok{=}\FloatTok{1.5}\NormalTok{,}
\NormalTok{            )}

            \CommentTok{\# Plot final condition}
            \FunctionTok{scatter!}\NormalTok{(}
\NormalTok{                ax,}
\NormalTok{                x\_data[}\KeywordTok{end}\NormalTok{],}
\NormalTok{                y\_data[}\KeywordTok{end}\NormalTok{],}
\NormalTok{                color}\OperatorTok{=}\NormalTok{ColorSchemes.glasbey\_hv\_n256[color\_idx],}
\NormalTok{                markersize}\OperatorTok{=}\FloatTok{12}\NormalTok{,}
\NormalTok{                marker}\OperatorTok{=:}\NormalTok{utriangle,}
\NormalTok{                strokecolor}\OperatorTok{=:}\NormalTok{black,}
\NormalTok{                strokewidth}\OperatorTok{=}\FloatTok{1.5}\NormalTok{,}
\NormalTok{            )}
        \ControlFlowTok{end}
    \ControlFlowTok{end}
\ControlFlowTok{end}

\CommentTok{\# Set global axis labels}
\FunctionTok{Label}\NormalTok{(}
\NormalTok{    fig[}\OperatorTok{:}\NormalTok{, }\KeywordTok{end}\NormalTok{, }\FunctionTok{Bottom}\NormalTok{()],}
    \StringTok{"phenotype 1"}\NormalTok{,}
\NormalTok{    fontsize}\OperatorTok{=}\FloatTok{14}\NormalTok{,}
\NormalTok{    padding}\OperatorTok{=}\NormalTok{(}\FloatTok{0}\NormalTok{, }\FloatTok{0}\NormalTok{, }\FloatTok{0}\NormalTok{, }\FloatTok{10}\NormalTok{),}
\NormalTok{)}

\FunctionTok{Label}\NormalTok{(}
\NormalTok{    fig[}\FloatTok{1}\NormalTok{, }\OperatorTok{:}\NormalTok{, }\FunctionTok{Left}\NormalTok{()],}
    \StringTok{"phenotype 2"}\NormalTok{,}
\NormalTok{    fontsize}\OperatorTok{=}\FloatTok{14}\NormalTok{,}
\NormalTok{    rotation}\OperatorTok{=}\ConstantTok{π} \OperatorTok{/} \FloatTok{2}\NormalTok{,}
\NormalTok{)}

\NormalTok{fig}
\end{Highlighting}
\end{Shaded}

\includegraphics[width=12.5in,height=12.5in]{evolutionary_dynamics_files/figure-pdf/cell-13-output-1.png}

In these visualizations, we can observe several key features of
evolutionary dynamics:

\begin{enumerate}
\def\labelenumi{\arabic{enumi}.}
\tightlist
\item
  \textbf{Fitness landscape influence}: Populations tend to climb
  fitness peaks (brighter regions).
\item
  \textbf{Mutational constraints}: Trajectories are influenced by the
  genotype-phenotype density map (white contour lines), avoiding regions
  of low density.
\item
  \textbf{Initial condition effects}: Different starting points can lead
  to different evolutionary outcomes.
\item
  \textbf{Stochasticity}: Even with the same starting point (replicate
  lines of the same color), stochasticity leads to different
  evolutionary paths.
\end{enumerate}

\subsection{Fitness Trajectory
Analysis}\label{fitness-trajectory-analysis}

Let's analyze how fitness changes over time during evolution:

\begin{Shaded}
\begin{Highlighting}[]
\BuiltInTok{Random}\NormalTok{.}\FunctionTok{seed!}\NormalTok{(}\FloatTok{42}\NormalTok{)}

\CommentTok{\# Define number of rows and columns}
\NormalTok{n\_rows }\OperatorTok{=} \FloatTok{4}
\NormalTok{n\_cols }\OperatorTok{=} \FloatTok{4}

\CommentTok{\# Define number of landscapes to display}
\NormalTok{n\_sample }\OperatorTok{=}\NormalTok{ n\_rows }\OperatorTok{*}\NormalTok{ n\_cols}

\CommentTok{\# Initialize figure}
\NormalTok{fig }\OperatorTok{=} \FunctionTok{Figure}\NormalTok{(size}\OperatorTok{=}\NormalTok{(}\FloatTok{600}\NormalTok{, }\FloatTok{600}\NormalTok{))}

\CommentTok{\# Add global grid layout}
\NormalTok{gl }\OperatorTok{=}\NormalTok{ fig[}\FloatTok{1}\NormalTok{, }\FloatTok{1}\NormalTok{] }\OperatorTok{=} \FunctionTok{GridLayout}\NormalTok{()}

\CommentTok{\# Loop over a sample of fitness landscapes}
\ControlFlowTok{for}\NormalTok{ (i, idx) }\KeywordTok{in} \FunctionTok{enumerate}\NormalTok{(}\FloatTok{1}\OperatorTok{:}\NormalTok{n\_sample)}
    \CommentTok{\# Define row and column indices}
\NormalTok{    row }\OperatorTok{=}\NormalTok{ (i }\OperatorTok{{-}} \FloatTok{1}\NormalTok{) }\OperatorTok{÷}\NormalTok{ n\_cols }\OperatorTok{+} \FloatTok{1}
\NormalTok{    col }\OperatorTok{=}\NormalTok{ (i }\OperatorTok{{-}} \FloatTok{1}\NormalTok{) }\OperatorTok{\%}\NormalTok{ n\_cols }\OperatorTok{+} \FloatTok{1}

    \CommentTok{\# Create subplot}
\NormalTok{    ax }\OperatorTok{=} \FunctionTok{Axis}\NormalTok{(}
\NormalTok{        gl[row, col],}
\NormalTok{        xlabel}\OperatorTok{=}\StringTok{"time (generations)"}\NormalTok{,}
\NormalTok{        ylabel}\OperatorTok{=}\StringTok{"fitness"}\NormalTok{,}
\NormalTok{        title}\OperatorTok{=}\StringTok{"landscape }\SpecialCharTok{$}\NormalTok{(idx)}\StringTok{"}\NormalTok{,}
\NormalTok{        titlesize}\OperatorTok{=}\FloatTok{12}\NormalTok{,}
\NormalTok{        yticklabelsvisible}\OperatorTok{=}\ConstantTok{false}\NormalTok{,}
\NormalTok{        xticklabelsvisible}\OperatorTok{=}\ConstantTok{false}\NormalTok{,}
\NormalTok{        xlabelsize}\OperatorTok{=}\FloatTok{10}\NormalTok{,}
\NormalTok{        ylabelsize}\OperatorTok{=}\FloatTok{10}\NormalTok{,}
\NormalTok{    )}

    \CommentTok{\# Plot fitness trajectories}
    \ControlFlowTok{for}\NormalTok{ lin }\KeywordTok{in}\NormalTok{ DD.}\FunctionTok{dims}\NormalTok{(x\_traj, }\OperatorTok{:}\NormalTok{lineage)}
        \ControlFlowTok{for}\NormalTok{ rep }\KeywordTok{in}\NormalTok{ DD.}\FunctionTok{dims}\NormalTok{(x\_traj, }\OperatorTok{:}\NormalTok{replicate)}
            \CommentTok{\# Extract fitness data}
\NormalTok{            fit\_data }\OperatorTok{=}\NormalTok{ x\_traj.fitness[}
\NormalTok{                fitness}\OperatorTok{=}\NormalTok{DD.}\FunctionTok{At}\NormalTok{(}\OperatorTok{:}\NormalTok{fitness),}
\NormalTok{                lineage}\OperatorTok{=}\NormalTok{lin,}
\NormalTok{                replicate}\OperatorTok{=}\NormalTok{rep,}
\NormalTok{                landscape}\OperatorTok{=}\NormalTok{idx,}
\NormalTok{                evo}\OperatorTok{=}\NormalTok{idx,}
\NormalTok{            ].data}

            \CommentTok{\# Plot trajectory}
\NormalTok{            color\_idx }\OperatorTok{=}\NormalTok{ (lin }\OperatorTok{{-}} \FloatTok{1}\NormalTok{) }\OperatorTok{*}\NormalTok{ n\_rep }\OperatorTok{+}\NormalTok{ rep}
            \FunctionTok{lines!}\NormalTok{(}
\NormalTok{                ax,}
                \FloatTok{0}\OperatorTok{:}\NormalTok{n\_steps,}
                \FunctionTok{vec}\NormalTok{(fit\_data),}
\NormalTok{                color}\OperatorTok{=}\NormalTok{ColorSchemes.glasbey\_hv\_n256[color\_idx],}
\NormalTok{                linewidth}\OperatorTok{=}\FloatTok{1}
\NormalTok{            )}
        \ControlFlowTok{end}
    \ControlFlowTok{end}
\ControlFlowTok{end}

\NormalTok{fig}
\end{Highlighting}
\end{Shaded}

\includegraphics[width=12.5in,height=12.5in]{evolutionary_dynamics_files/figure-pdf/cell-14-output-1.png}

\subsection{Conclusion}\label{conclusion}

In this notebook, we've simulated evolutionary dynamics in a 2D
phenotypic space under various fitness landscapes and a fixed mutational
landscape. We've used the Metropolis-Hastings algorithm to model
evolution as a biased random walk influenced by both fitness and
mutational constraints.

Key insights from our simulations:

\begin{enumerate}
\def\labelenumi{\arabic{enumi}.}
\tightlist
\item
  \textbf{Fitness-mutation interplay}: Evolution is shaped by both the
  fitness landscape (selection) and the mutational landscape (genetic
  constraints).
\item
  \textbf{Multiple evolutionary outcomes}: Different initial conditions
  can lead to different evolutionary endpoints, even in the same fitness
  landscape.
\item
  \textbf{Stochastic effects}: Even with identical starting points,
  stochasticity in the evolutionary process leads to variation in
  outcomes.
\item
  \textbf{Landscape navigation}: Populations navigate the phenotypic
  space by climbing fitness gradients while being constrained by the
  mutational landscape.
\end{enumerate}

This approach allows us to explore how populations might adapt to
various environments given underlying genetic constraints, providing
insights into the predictability and paths of evolution.



\end{document}
